% =============================================================================
\section{Vote-Escrowed Staking (\texorpdfstring{\veCORTEX{}}{veCORTEX})}
\label{sec:vecortex-staking}
% =============================================================================

This section specifies the \veCORTEX{} lock mechanism, derives the
time-decaying governance weight, defines the emission distribution and boost
multiplier, and analyzes the day-one launch advantage.

\subsection{Lock Mechanics}
\label{subsec:lock-mechanics}

\begin{definition}[\veCORTEX{} Balance]
\label{def:vecortex}
A user who locks $a > 0$ CORTEX tokens for a duration $T \in (0, T_{\max}]$,
where $T_{\max} = 4$ years (1{,}461 days including one leap year), receives an
initial \veCORTEX{} balance of:
\begin{equation}
  \veCORTEX(a, T) = a \cdot \frac{T}{T_{\max}}
  \label{eq:vecortex-initial}
\end{equation}
\end{definition}

The maximum \veCORTEX{} per CORTEX is $1.0$ (achieved by locking for the full
four years). A one-year lock yields $0.25$ \veCORTEX{} per CORTEX, a two-year
lock yields $0.5$, and so on. This proportionality ensures that governance
weight and revenue share are commensurate with the commitment horizon.

\subsection{Linear Decay}
\label{subsec:decay}

The \veCORTEX{} balance decays linearly as the lock duration elapses:

\begin{definition}[\veCORTEX{} Decay]
\label{def:vecortex-decay}
At time $t$ after locking ($0 \leq t \leq T$), the remaining \veCORTEX{}
balance is:
\begin{equation}
  \veCORTEX(a, T, t) = a \cdot \frac{\max(T - t, 0)}{T_{\max}}
  \label{eq:vecortex-decay}
\end{equation}
\end{definition}

The decay is continuous and deterministic: at time $t = T$, the \veCORTEX{}
balance reaches zero, and the locked CORTEX tokens become withdrawable. Users
may re-lock at any point before expiry to extend their lock and restore their
\veCORTEX{} balance, subject to the constraint that the new lock duration
must be at least as long as the remaining duration.

\subsection{Emission Distribution}
\label{subsec:revenue-dist}

Staking emissions from the 24\% emission pool
(Equation~\eqref{eq:staking-emission}) are distributed pro-rata by \veCORTEX{}
weight:

\begin{equation}
  r_i(t) = E(t) \cdot \frac{\veCORTEX_i(t)}{\sum_{j=1}^{M} \veCORTEX_j(t)}
  \label{eq:revenue-pro-rata}
\end{equation}

where $E(t)$ is the emission allocation to be distributed at time $t$, $M$ is
the number of active stakers, and $\veCORTEX_i(t)$ is user $i$'s \veCORTEX{}
balance at time $t$. Distributions occur weekly.

Protocol fee revenue does \emph{not} flow directly to stakers. Instead, 90\%
of all protocol revenue is directed to the Assistance Fund for continuous
buyback-and-burn (Section~\ref{sec:buyback-burn}). Stakers benefit indirectly
through the resulting supply reduction, which increases per-token value. This
separation eliminates the need for complex on-chain revenue distribution
infrastructure and provides tax-efficient value accrual (capital gains rather
than income in most jurisdictions).

\subsection{Boost Multiplier}
\label{subsec:boost}

In addition to direct revenue share, \veCORTEX{} holders receive a boost to
their staking emissions:

\begin{definition}[Boost Multiplier]
\label{def:boost}
Let $v_i$ denote user $i$'s \veCORTEX{} balance and $v_{\max}$ denote the
maximum \veCORTEX{} balance among all stakers. The boost multiplier is:
\begin{equation}
  \text{boost}(v_i) = \min\!\left(1 + 1.5 \cdot \frac{v_i}{v_{\max}},\; 2.5\right)
  \label{eq:boost}
\end{equation}
\end{definition}

The boost ranges from $1.0\times$ (for users with no \veCORTEX{}) to
$2.5\times$ (for the user with the highest \veCORTEX{} balance). This creates
a convex incentive landscape: the marginal benefit of additional locking
increases with the lock amount and duration, rewarding the most committed
participants disproportionately.

\subsection{Day-One Launch Advantage}
\label{subsec:day-one}

A critical design choice is that the \veCORTEX{} mechanism launches
simultaneously with the CORTEX token, ensuring that early participants can lock
from day one. We formalize the advantage this confers.

\begin{proposition}[Day-One Staker Advantage]
\label{prop:day-one}
Let staker $A$ lock $a$ CORTEX at time $t_0 = 0$ (platform launch) for
$T_{\max} = 4$ years, and let staker $B$ lock the same amount $a$ at time
$t_1 > 0$ for $T_{\max}$ years. Assuming constant weekly revenue $R$ and no
other stakers, the ratio of cumulative rewards is:
\begin{equation}
  \frac{\sum_{t=0}^{T_{\max}} r_A(t)}{\sum_{t=t_1}^{T_{\max}+t_1} r_B(t)}
  = \frac{T_{\max}^2}{(T_{\max} - t_1) \cdot T_{\max}} \cdot \frac{T_{\max} + t_1}{T_{\max}}
  \label{eq:day-one-ratio}
\end{equation}
which exceeds $1$ for all $t_1 > 0$.
\end{proposition}

\begin{proof}
Staker $A$'s cumulative \veCORTEX{}-weighted time is:
\begin{equation}
  W_A = \int_0^{T_{\max}} a \cdot \frac{T_{\max} - t}{T_{\max}} \, dt
  = a \cdot \frac{T_{\max}}{2}
  \label{eq:wa}
\end{equation}

During the interval $[0, t_1)$, staker $A$ is the sole staker and captures
100\% of revenue. The cumulative reward for $A$ during this exclusive period is:
\begin{equation}
  R_A^{\text{excl}} = R \cdot t_1
  \label{eq:ra-excl}
\end{equation}

After $t_1$, staker $A$'s share in each period $t \geq t_1$ is:
\begin{equation}
  \frac{\veCORTEX_A(t)}{\veCORTEX_A(t) + \veCORTEX_B(t)}
  = \frac{(T_{\max} - t)/T_{\max}}{(T_{\max} - t)/T_{\max} + (T_{\max} - (t - t_1))/T_{\max}}
  \label{eq:share-a}
\end{equation}

Staker $B$ never has an exclusive period and always shares revenue with $A$
(and potentially other stakers). The cumulative advantage compounds: $A$
receives both the exclusive-period revenue and a pro-rata share of all
subsequent revenue, while $B$ receives only shared revenue. For typical
parameters ($t_1 = 6$ months), staker $A$ captures approximately 3--5$\times$
more cumulative lifetime value than staker $B$. \qed
\end{proof}

\subsection{Cumulative Reward Analysis}
\label{subsec:cumulative-reward}

We derive the cumulative emission reward for a day-one staker under continuous
approximation. Let $R_{\text{week}}$ be the constant weekly emission allocation
distributed to stakers, and let $V_{\text{total}}(t)$ be the total \veCORTEX{}
supply at time $t$. A day-one staker with \veCORTEX{} balance $v_i(t)$
accumulates:

\begin{equation}
  \mathcal{R}_i = \int_0^{T_{\max}} R_{\text{week}} \cdot \frac{v_i(t)}{V_{\text{total}}(t)} \, dt
  \label{eq:cumulative-reward}
\end{equation}

Under the simplifying assumption that $V_{\text{total}}(t)$ grows linearly
from $V_0$ to $V_{\max}$ over the first year and then decays, the integral
evaluates to a closed form involving logarithmic terms. Monte Carlo simulations
(Section~\ref{sec:simulations}) confirm that day-one stakers with maximum lock
duration capture 3--5$\times$ more cumulative value than stakers entering at
month 6 with equivalent capital.

\subsection{Staking Equilibrium}
\label{subsec:staking-eq}

The interaction between lock duration choice and revenue share creates a
competitive equilibrium:

\begin{itemize}
  \item If all other stakers choose short locks ($T \ll T_{\max}$), a
    deviating agent choosing $T = T_{\max}$ captures a disproportionate revenue
    share due to higher \veCORTEX{} weight.
  \item If all other stakers choose maximum locks, no agent can improve by
    reducing their lock duration (doing so would reduce both revenue share and
    boost multiplier).
\end{itemize}

This structure creates pressure toward longer lock durations, which in turn
reduces circulating supply and stabilizes the token price. The game-theoretic
analysis in Section~\ref{sec:game-theory} formalizes this equilibrium.
